Put
\[
 u:=\frac{d^2}{n^2\log^2(en)},\qquad h:=\sigma^2+\frac d{n^2}.
\tag{1}
\]
Theorem~\(\mathrm{thm:robust\text{-}upper\text{-}construction\text{-}resolution}\) constructs two total, clipped, data-driven estimators and gives
\[
 \sup_{P\in\mathcal P_{d,\epsilon,M,\sigma}}
 E_{Q_P^{(n)}}[(T_n^{\mathrm{poly}}-\tau(P))^2]
 \le C_\epsilon M^2\{n^{-1}+\min(1,u)\}
\tag{2}
\]
and
\[
 \sup_{P\in\mathcal P_{d,\epsilon,M,\sigma}}
 E_{Q_P^{(n)}}[(T_n^{\mathrm{col}}-\tau(P))^2]
 \le C_\epsilon M^2(n^{-1}+h).
\tag{3}
\]
Take the smaller of the two dimension constants and the larger of the two risk constants in that theorem.

If \(h\le\min\{1,u\}\), Definition~\(\mathrm{def:total\text{-}estimator}\) selects \(T_n^{\mathrm{col}}\), and the remainder in (3) is
\[
 h=\min\{1,u,h\}.
\tag{4}
\]
If \(u<\min\{1,h\}\), the selector chooses \(T_n^{\mathrm{poly}}\).  This inequality entails \(u<1\), so the remainder in (2) is
\[
 u=\min\{1,u,h\}.
\tag{5}
\]
Suppose neither selector inequality holds.  Then \(\min\{1,u,h\}=1\).  Indeed, if that minimum were smaller than one, either \(h\le u\), which would trigger the first case, or \(u<h\), which would trigger the second case.  In the remaining case the definition selects zero.  Lemma~\(\mathrm{lem:scale\text{-}sanity}\) gives \(\tau(P)\in[-M,M]\), and hence
\[
 (0-\tau(P))^2\le M^2
 \le M^2\{n^{-1}+\min(1,u,h)\}.
\tag{6}
\]
Equations (2)--(6) cover every tie and boundary case and prove
\[
 \sup_{P\in\mathcal P_{d,\epsilon,M,\sigma}}
 E_{Q_P^{(n)}}[(\widehat\tau_n^\star-\tau(P))^2]
 \le C_\epsilon M^2\left\{\frac1n+
 \min\left(1,\frac{d^2}{n^2\log^2(en)},\sigma^2+\frac d{n^2}\right)\right\}.
\tag{7}
\]

For the published comparison, Lemma~\(\mathrm{lem:zeng\text{-}binary\text{-}collision\text{-}upper}\), followed by the identity \(E[(Z-\theta)^2]=\operatorname{Var}(Z)+\{E(Z)-\theta\}^2\), gives the binary collision benchmark with heterogeneity-dependent remainder of order \(h\) on the present dimension range; its exponentially small bias term is absorbed by \(n^{-1}+d/n^2\).  Since
\[
 \min\{1,u,h\}\le h,
\tag{8}
\]
the selector remainder is never larger in order.  If \(d=o(n\log(en))\), then \(u=o(1)\).  If also \(u=o(h)\), eventually \(u<\min\{1,h\}\), so the selector remainder is \(u=o(h)\).  This proves the strict-order comparison.  The construction uses known \(M\) and \(\sigma\); it makes no inference or unknown-parameter adaptation claim.